\documentclass[11pt]{article}
\usepackage[T1]{fontenc}
\usepackage{amsmath,amssymb,amsthm}
\usepackage[margin=1in]{geometry}
\usepackage{booktabs}
\usepackage{microtype}
\usepackage{url}
\usepackage[hidelinks,
  pdftitle={The least Type 1 cubic graphs without equitable 4-total colorings have order 16},
  pdfauthor={Xiaojian Zeng}]{hyperref}

\newtheorem{theorem}{Theorem}
\newtheorem{lemma}[theorem]{Lemma}

\title{The least Type~1 cubic graphs without equitable
  \(4\)-total colorings have order \(16\)}
\author{Xiaojian Zeng\\
PBC School of Finance, Tsinghua University, 100190\\
\texttt{xiaojianzeng@foxmail.com}}
\date{}

\begin{document}
\maketitle

\begin{abstract}
Adauto, de Figueiredo, Sasaki, and Schneider asked whether every Type~1
cubic graph of order less than \(20\) has an equitable proper
\(4\)-total coloring.  We give a negative answer by an exact finite
classification.  The minimum counterexample order is \(16\), and there
are exactly \(16\) isomorphism classes at that order.  A canonical census
contains all \(45{,}982\) connected simple cubic graphs of orders
\(4,6,\ldots,18\).  Every graph through order \(16\) is covered by a
directly checked equitable coloring, a checked proof that no proper
\(4\)-total coloring exists, or---for the \(16\) counterexamples---a
directly checked proper coloring together with checked unsatisfiability
certificates for equitable colorability.  Census completeness is checked
independently by exhaustive canonization and exact orbit sums against a
labeled-graph recurrence.  Thus the previously known order-\(20\) graph
is not minimum.  No classification at order \(18\), or claim at larger
orders, is made.
\end{abstract}

\section{Context and result}

All graphs in this article are finite, simple, and undirected.  A total
coloring assigns colors to vertices and edges so that adjacent vertices,
adjacent edges, and incident vertex--edge pairs receive distinct colors.
It is equitable when any two color classes differ in cardinality by at
most one.  A cubic graph is \emph{Type~1} when its total chromatic number
is \(4\).

Dantas et al.\ exhibited Type~1 cubic graphs whose equitable total
chromatic number is \(5\) and proved that equitable \(4\)-total
colorability is NP-complete even for bipartite cubic graphs
\cite{DantasEtAl2016}.  One of their examples is the order-\(20\) graph
\(R\).  Adauto et al.\ subsequently distinguished the existence of an
equitable coloring from the stronger requirement that every
\(4\)-total coloring be equitable.  Their Question~7.1 asks whether every
Type~1 cubic graph below order \(20\) has at least one equitable
\(4\)-total coloring, and suggests a complete catalogue computation by
integer programming or SAT \cite{AdautoEtAl2026}.

The contribution here is an exact negative resolution of that finite
question, including minimum order and the complete least-order
classification.  The computation uses canonical graph representatives,
explicit positive witnesses, proof-carrying SAT, and a census-completeness
argument independent of the generator.  The practice of coupling
canonical graph6 representatives with a definition-level independent
verifier also appears in a recent exact finite graph classification by
Zeng \cite{Zeng2026Whiskered}; the algebraic property treated there is
unrelated to total coloring.

\begin{theorem}\label{thm:main}
The assertion that every connected Type~1 cubic graph \(G\) with
\(|V(G)|<20\) admits an equitable proper \(4\)-total coloring is false.
The minimum order of a counterexample is \(16\), and exactly \(16\)
isomorphism classes are counterexamples at that order.
\end{theorem}

In canonical census order, their graph6 codes are
\begin{center}
\small
\begin{tabular}{ll}
\verb|Os?GOCF@O??ACBCBL?B?O| & \verb|Ot?GOCA?WB@?ODOBC_?g@|\\
\verb|Os?GO?C?_MAECKOOCGJO?| & \verb|Os?GOKE??@`?CDCBKABA?|\\
\verb|Os?GGOB@OO?OGBCBGEB_?| & \verb|Ot?GOKFA??_A?B?BK?Q_A|\\
\verb|Ov?GW?@?OM?E?KA?G?h?E| & \verb|Os?GOGA?AOgoGWGWAA?_`|\\
\verb|Os??WWK?@CaGCGOCCAI?K| & \verb|Ov?GO???WB?k?c?a?_ZO?|\\
\verb|Os??WWA?W@?ACBCBL?BG?| & \verb|Os??WOB?OW?ACBCDKCBG?|\\
\verb|Os??WOC?@@cEOSQOAO?g@| & \verb|Os???[AAGH@_C@GEIOB?C|\\
\verb|Os?GG???WBD_CcCaGKBO?| & \verb|Os?GW?@?OMA?E@OKCEB?G|\\
\end{tabular}
\end{center}

\section{Coloring formulas}

Let \(G\) be cubic of order \(n\).  It has \(3n/2\) edges, so a total
coloring colors \(5n/2\) objects.  Write \(A(G)\) for the existence of a
proper \(4\)-total coloring and \(E(G)\) for the existence of an equitable
one.  Since every cubic graph has total chromatic number at least \(4\),
Type~1 is equivalent to \(A(G)\).  The question is therefore whether
\(A(G)\Rightarrow E(G)\) for every eligible \(G\).

The following identity, also given as the decomposition lemma of
Adauto et al.\ \cite{AdautoEtAl2026}, is useful both mathematically and as
the basis of a second encoding.

\begin{lemma}[four-color star identity]\label{lem:star}
Fix a proper \(4\)-total coloring of a cubic graph of order \(n\).  For a
color \(c\), let \(S_c\) be the vertices of color \(c\), and let \(M_c\)
be the edges of color \(c\).  Then
\[
  n=|S_c|+2|M_c|,
  \qquad
  |S_c|+|M_c|=\frac{n+|S_c|}{2}.
\]
In particular, at order \(16\) the coloring is equitable if and only if
\(|S_c|=4\) for every color \(c\).
\end{lemma}

\begin{proof}
At each vertex, that vertex and its three incident edges are pairwise
conflicting, hence use all four colors.  Counting occurrences of a fixed
color over the \(n\) vertex-stars counts a vertex in \(S_c\) once and an
edge in \(M_c\) at both endpoints.  This proves the first identity, and
the second follows by elimination.  At order \(16\) there are \(40\)
objects, so equitability requires four classes of size \(10\); the second
identity makes this equivalent to \(|S_c|=4\) for all \(c\).
\end{proof}

For the primary CNF, order the objects as vertices \(0,\ldots,n-1\)
followed by lexicographically ordered edges, and introduce \(x_{a,c}\)
for every object \(a\) and color \(c\in\{0,1,2,3\}\).  An at-least-one
clause and the six pairwise negative clauses impose exactly one color on
each object.  For every conflicting pair \(a,b\) and every \(c\), add
  \(\lnot x_{a,c}\lor\lnot x_{b,c}\).  Finally, \(x_{0,0}\) fixes a harmless
color symmetry.  The resulting formula is satisfiable exactly when
\(A(G)\) holds.

For \(E(G)\), the primary encoding additionally bounds every color-class
sum between \(\lfloor(5n/2)/4\rfloor\) and
\(\lceil(5n/2)/4\rceil\), using sequential-counter clauses
\cite{Sinz2005} as implemented by PySAT \cite{IgnatievEtAl2018}.  These
bounds are equivalent to equitability because every object has exactly
one color.  All retained DIMACS files were regenerated from their graph6
rows and compared clause for clause.

The decisive order-\(16\) failures have a second, independently written
encoding.  It uses only the \(160\) variables \(x_{a,c}\), the direct
proper-total-coloring clauses above, and, for every color \(c\) and every
five-element set \(U\subseteq V(G)\), the clause
\[
  \mathop{\vee}\limits_{v\in U}\lnot x_{v,c}.
\]
Thus at most four vertices receive any one color.  The \(16\) exactly-one
vertex assignments force equality for all four colors, and
Lemma~\ref{lem:star} proves that this is exactly equitable
\(4\)-total colorability.  Each such formula has \(18{,}233\) clauses.

Unsatisfiability is carried by DRAT derivations checked against the exact
DIMACS inputs.  DRAT and its checker are described by Wetzler, Heule, and
Hunt \cite{WetzlerEtAl2014}.  Solver exit status is never used as a
certificate: a negative classification is admitted only after a retained
derivation is replayed successfully.

\section{Complete census}

Simplicity and the handshaking lemma restrict the orders below \(20\) to
\(4,6,\ldots,18\).  Representatives were generated with the
\texttt{geng} program in the nauty suite, using
\begin{center}
\texttt{geng -cld3D3 n (3*n/2)}.
\end{center}
The options request connected graphs, canonical labels, and minimum and
maximum degree three.  The nauty refinement--individualization framework
is described by McKay and Piperno \cite{McKayPiperno2014}.  Every graph6
row was decoded and checked for order, simplicity, cubicity, and
connectedness.

Completeness does not rest on the generator's status output.  A separate
C program starts from the one-cell vertex partition, repeatedly refines
vertices by their neighbor-count vectors into the current cells, and, if
the stable partition is not discrete, individualizes every vertex in an
isomorphism-invariantly selected nonsingleton cell.  It performs no branch
pruning.  Consequently isomorphic inputs have the same multiset of leaf
adjacency strings.  The least string is a canonical form, and the number
of leaves attaining it is \(|\operatorname{Aut}(G)|\).  The census has
distinct canonical forms at every order.

An independent recurrence counts the labeled graphs.  Let
\(F(c_1,c_2,c_3)\) count simple completions of a fixed labeled residual
degree sequence with \(c_i\) vertices of residual degree \(i\).  Remove a
fixed vertex of maximum residual degree \(d\); choosing \(k_i\) of the
remaining degree-\(i\) vertices as its neighbors, where
\(k_1+k_2+k_3=d\), lowers each selected residual degree by one.  Summing
the corresponding products of binomial coefficients gives \(F\), with
\(F(0,0,0)=1\).  Hence \(L_n=F(0,0,n)\) counts all labeled simple cubic
graphs.  If \(C_n\) counts the connected ones and label \(1\) lies in a
component of order \(k\), then
\[
 C_n=L_n-
 \sum_{\{k\in 2\mathbb{Z}:4\leq k<n\}}
 \binom{n-1}{k-1}C_kL_{n-k}.
\]
This counts each graph once, by the unique component containing label
\(1\).

The resulting values and the census orbit sums agree exactly:
\begin{center}
\begin{tabular}{r@{\qquad}r@{\qquad}r}
\toprule
\(n\) & isomorphism classes & \(C_n=\sum_G n!/|\operatorname{Aut}(G)|\)\\
\midrule
4  & 1      & 1\\
6  & 2      & 70\\
8  & 5      & 19,320\\
10 & 19     & 11,166,120\\
12 & 85     & 11,543,439,600\\
14 & 509    & 19,491,385,914,000\\
16 & 4,060  & 50,233,275,604,512,000\\
18 & 41,301 & 187,663,723,374,359,232,000\\
\bottomrule
\end{tabular}
\end{center}
The class counts total \(45{,}982\).  Pairwise nonisomorphism and the
orbit-sum equality prove completeness: any omitted isomorphism class
would contribute the positive integer \(n!/|\operatorname{Aut}(G)|\) and
make the displayed orbit sum smaller than the independent labeled count.

\section{Classification and proof of Theorem~\ref{thm:main}}

The exhaustive join through order \(16\) is as follows.  Here
``\(A\)-UNSAT'' means a checked DRAT refutation of the proper-coloring
formula, and the last column combines a directly checked proper coloring
with a checked refutation of the equitable formula.

\begin{center}
\begin{tabular}{r@{\qquad}r@{\qquad}r@{\qquad}r}
\toprule
order & equitable witness & \(A\)-UNSAT & proper witness + \(E\)-UNSAT\\
\midrule
4  & 0    & 1   & 0\\
6  & 1    & 1   & 0\\
8  & 4    & 1   & 0\\
10 & 9    & 10  & 0\\
12 & 60   & 25  & 0\\
14 & 461  & 48  & 0\\
16 & 3,817& 227 & 16\\
\bottomrule
\end{tabular}
\end{center}

There are \(4{,}352\) equitable witness arrays and \(16\) further proper
witness arrays.  A verifier reconstructs the ordered vertex and edge
objects from the graph6 row and checks every vertex, edge, and incidence
conflict and every asserted class size directly.  Each of the \(16\)
proper witnesses has sorted total-class profile \((9,10,10,11)\), which
establishes Type~1 but is not used to infer equitable nonexistence.

For the remaining rows, \(313\) primary \(A\)-formulas and \(16\) primary
\(E\)-formulas have replayed DRAT refutations.  All \(16\) decisive
\(E\)-refutations are additionally replayed for the independent
\(160\)-variable encoding derived from Lemma~\ref{lem:star}.  Thus each
graph below order \(16\) either has an equitable witness or is not Type~1.
At order \(16\), precisely the \(16\) listed codes have a proper witness
and no equitable coloring.  The complete census then proves both
minimality and the exact least-order classification, completing the proof
of Theorem~\ref{thm:main}. \qed

Order \(18\) was enumerated to complete the specified below-\(20\)
census, but it was not color-classified: after a counterexample at order
\(16\) and a complete classification at all lower orders, such a
classification is unnecessary for the theorem.  The result does not
characterize all larger counterexamples.

\section{Reproducibility and limitations}

The accompanying artifact has one verification entry point,
\path{evidence/scripts/final_verify.py}.  It rebuilds the retained
generator and two checkers, reproduces every census file byte for byte,
reruns the independent canonization and labeled-count test, directly
checks all \(4{,}368\) positive coloring arrays, regenerates and binds the
\(329\) primary CNFs, and replays the \(329\) primary and \(16\) secondary
DRAT derivations.  The recorded environment uses Python~3.10.12,
NetworkX~3.4.2, PySAT~1.8.dev24, and Zig~0.13.0.  Census generation uses
no random option; the deterministic relabeling test of the independent
canonizer records seed \(6802\).  Certificate generation requested no
randomization.

The retained nauty generator is tied to version~2.8.9 by its object
provenance, and can be relinked byte identically, but its full source
archive is not part of the artifact.  This is a provenance limitation of
the generator.  It is not a completeness assumption: the separately
implemented canonizer and labeled recurrence establish the census
independently.  Likewise, the second equitable encoding avoids reliance
on the primary cardinality library for every decisive counterexample.

\paragraph{Computational and AI assistance.}
Python and C programs performed graph generation, exact enumeration,
SAT encoding, witness checking, and certificate replay.  OpenAI Codex
assisted with the research workflow, verification code, and manuscript
drafting.  No solver status or AI-generated assertion was treated as a
mathematical proof; the conclusions above rest on explicit witnesses,
definition-level checks, exact integer identities, and replayed proof
certificates.

\bibliographystyle{plainurl}
\bibliography{references}

\end{document}
